% Generated draft: requires mathematical and source review.
\documentclass[11pt]{article}
\usepackage[margin=25mm]{geometry}
\usepackage{fontspec}
\setmainfont{Latin Modern Roman}
\usepackage{amsmath,amssymb,mathtools,unicode-math}
\setmathfont{Latin Modern Math}
\usepackage{enumitem,needspace,xurl,hyperref}
\setlength{\parindent}{0pt}
\setlength{\parskip}{5pt}
\setlength{\emergencystretch}{3em}
\newcommand{\N}{\mathbb{N}}
\newcommand{\Z}{\mathbb{Z}}
\newcommand{\Q}{\mathbb{Q}}
\newcommand{\R}{\mathbb{R}}
\newcommand{\C}{\mathbb{C}}
\newcommand{\F}{\mathbb{F}}
\newcommand{\A}{\mathbb{A}}
\newcommand{\PP}{\mathbb{P}}
\newcommand{\E}{\mathbb{E}}
\newcommand{\eps}{\varepsilon}
\newcommand{\dd}{\,\mathrm{d}}
\begin{document}
\section*{UnsolvedMath ID: 30006601}
\textbf{Source problem code:} catalogue entry
\subsection{Definitions and mathematical statement}
The intended assertion concerns every finite-dimensional algebra \(\Lambda\). The supplied record does not specify its ground field, additional hypotheses, or the category \(K_\Lambda\). Write \(I_\Lambda=\operatorname{ind}K_\Lambda\) for the index collection appearing in the source; its objects and the convention for choosing indecomposable representatives remain unspecified.

The source uses polynomial and Laurent polynomial coordinate rings over \(\C\), provisionally written as
\[
R_\Lambda^0=\C[u_X\mid X\in I_\Lambda],\qquad
R_\Lambda^\times=\C[u_X^{\pm1}\mid X\in I_\Lambda].
\]
Here \(\widehat F_X\) is a provisional transcription of the universal \(F\)-polynomial symbol in the damaged excerpt, and \(q_X=0\) provisionally denotes an indexed \(u\)-equation. Neither family is defined in the supplied record. Subject to confirming these generators and their indexing, the affine definitions can be transcribed as
\[
\widetilde{\mathcal M}_\Lambda
=\operatorname{Spec}\left(R_\Lambda^0/
\langle\widehat F_X-1\mid X\in I_\Lambda\rangle\right),
\qquad
\widetilde{\mathcal U}_\Lambda
=\operatorname{Spec}\left(R_\Lambda^0/
\langle q_X\mid X\in I_\Lambda\rangle\right).
\]
The source also defines untilded spaces \(\mathcal M_\Lambda\) and \(\mathcal U_\Lambda\) using Laurent coordinates. Its transcription does not reliably distinguish the ordinary and universal \(F\)-polynomial symbols or all coordinate decorations in those definitions. Consequently, the precise Laurent defining ideals, including whether they are extensions of the displayed polynomial ideals, require confirmation.

The requested conjecture is that, for every finite-dimensional algebra \(\Lambda\) covered by the intended definitions,
\[
\widetilde{\mathcal M}_\Lambda=\widetilde{\mathcal U}_\Lambda,
\qquad
\mathcal M_\Lambda=\mathcal U_\Lambda.
\]
The meaning of equality also requires confirmation. If these are equalities of closed subschemes in the indicated common ambient affine space and torus, they require equality of the corresponding defining ideals; equality of radicals would establish only equality of the associated reduced closed subschemes. The supplied information does not resolve this distinction or supply the missing generators, so this is a qualified statement schema rather than a complete intrinsic formulation.
\subsection{Short English statement}
For every finite-dimensional algebra in the intended scope, the conjecture asks whether the configuration spaces defined by universal \(F\)-polynomials coincide with those defined by the \(u\)-equations, both in polynomial and Laurent coordinates. The supplied excerpt omits essential definitions and has damaged notation, so the precise hypotheses, equations, and meaning of equality require human review.
\subsection{Sources}
\begin{itemize}
\item Supplied local UnsolvedMath registry record, permanent ID \(30006601\), title: Universal \(F\)-Polynomial Configuration-Space Identities; exact source problem code: \texttt{catalogue entry}. The record supplies the conjectural identities and the damaged source excerpt.
\item Representation Theory of Quivers and Finite-Dimensional Algebras (2026), Oberwolfach Reports, DOI 10.4171/owr/2026/6, pp.\ 345--432. Supplied source URL and literature reference labelled Oberwolfach Report 6/2026: \url{https://doi.org/10.4171/owr/2026/6}. The supplied assessment identifies the universal \(F\)-polynomial equality as Conjecture 7. No individual proposer is specified.
\item Klasz--Marczinzik--Thomas, Auslander regular algebras, supplied literature reference: \url{https://arxiv.org/abs/2501.09447}. Its precise role in the conjecture is not explained in the supplied record.
\item Literature assessment checked at \texttt{2026-08-22}: this is supplied dated curation, not an independently verified assessment. It labels the problem open and reports that the conjecture is known for Dynkin path algebras and several other families, without reporting a general proof. The background expressly requests independent verification of time-sensitive status and cited claims.
\item The registry labels the statement \texttt{corrected\_verified} and says it was reconstructed and checked against the official report. These are supplied source-verification claims; they have not been independently verified in this draft.
\end{itemize}
\end{document}
